% \WebLikePDF in the document's config (test-likepdf.cfg), which sets the
% text as the PDF does. The checks are in test/run.mjs, those of the sizes
% in test/page.mjs, and those of the space around center and verbatim, a
% list's plus \parskip, in test/visual.mjs. A rule as wide as the text
% after a paragraph's indent stays in the column of a phone's screen (the
% check that the page fits in test/visual.mjs).
\documentclass[11pt]{article}
\usepackage{amsmath}
\setlength{\parskip}{3pt}
\begin{document}
\tableofcontents
\section{Paragraphs}
The first paragraph after a heading isn't indented in the PDF. Its lines are
long enough to wrap, so that the justified text and its hyphenation show on a
wide screen: notwithstanding the incomprehensibility of characteristically
uncharacteristic words, the text reaches the right margin.\footnote{A note,
which stays ragged right beside the text.}

A second paragraph is indented, with \verb|\parskip| above it, and has inline math
such as $\sum_{i=1}^n x_i^2 \le \bigl(\sum_{i=1}^n |x_i|\bigr)^2$ in a line
that has to break around it somewhere.
\[ \int_0^1 f(x)\,dx = F(1) - F(0). \]
A paragraph that goes on after the display.
\begin{itemize}
\item An item long enough to wrap onto a second line, whose text is also
  justified on a wide screen, as LaTeX justifies a list's items.
\end{itemize}
A paragraph before a center.
\begin{center}
    Centered.
\end{center}
A paragraph before a verbatim.
\begin{verbatim}
verbatim text
\end{verbatim}
A paragraph with a footnote.\footnote{A footnote.}
\begin{center}
    Centered after a footnote.
\end{center}
A paragraph after the blocks.

\rule{\textwidth}{0.4pt}, an indented rule as wide as the text, in the column.
\end{document}
